\SourceLeaf{431}%
\[
\begin{aligned}
\frac1r&=A\cos\varphi^{2}+2B\cos\varphi\sin\varphi+C\sin\varphi^{2}\\
&=\frac{A+C}{2}+\frac{A-C}{2}\cos2\varphi+B\sin2\varphi.
\end{aligned}
\]
Macht man
\[
\begin{aligned}
\frac{A-C}{2}&=E\cos2\theta\\
B&=E\sin2\theta,
\end{aligned}
\]
wo man annehmen kann, dass $E$ dasselbe Zeichen wie $\frac{A-C}{2}$ hat, so wird
\[
\frac1r=\tfrac12(A+C)+E\cos2(\varphi-\theta).
\]
Offenbar bedeutet $\varphi$ den Winkel zwischen der schneidenden Ebene und einer andern durch die Normale und diejenige Tangente gehenden, der die Richtung $M$ entspricht. Offenbar hat also $\frac1r$ seinen grössten (absoluten) oder $r$ seinen kleinsten Werth, wenn $\varphi=\theta$, und $\frac1r$ den kleinsten, wenn $\varphi=\theta+90^{0}$ wird; die grösste und kleinste Krümmung finden also in zweien einander unter rechten Winkeln schneidenden Ebenen statt. Diese äussersten Werthe für $\frac1r$ sind folglich
\[
\tfrac12(A+C)\pm\BareRoot{\left\{\left(\frac{A-C}{2}\right)^{2}+BB\right\}},
\]
ihre Summe $=A+C$ und ihr Product $=AC-BB$ oder das Product der beiden äussersten Krümmungshalbmesser
\[
=\frac{1}{AC-BB}.
\]
Dieses Product, welches von grosser Wichtigkeit ist, verdient genauer entwickelt zu werden. Wir finden nemlich aus obigen Formeln
\[
AC-BB=(\alpha\beta'-\beta\alpha')^{2}(TV-UU)ZZ.
\]
Nach der dritten Formel in [Lehrsatz] $6$, Art. $7$, folgert man aber leicht, dass
\[
\alpha\beta'-\beta\alpha'=\pm Z,
\]
also
\[
AC-BB=Z^{4}(TV-UU).
\]
Übrigens ist nach Art. $8$
